\section{Second variation formula and Bonnet-Myers theorem}
Let $M$ be a Riemannian manifold, and $c:(-\epsilon, \epsilon)\times [0,\ell] \to M$ a $C^{\infty}$ map.
Set $c_s(t)=c^t(s):=c(s,t)$. Assume $c_0$ is a non-trivial geodesic.
Let $X:=dc(\frac{\partial}{\partial s}), \quad T:= dc(\frac{\partial}{\partial t}), \quad X^{\perp}:= X- \langle X, \frac{T}{\|T\|}\rangle \frac{T}{\|T\|}$.
$X^{\perp}$ is the component of $X$ orthogonal to $T$. Since $c_0$ is non-trivial, $T(0,t) = c_0`(t) \neq 0$.
By continuity, $T(s,t) \neq 0$ for sufficiently small $s$. Thus, by replacing $\epsilon >0$ if necessary, assume these are vector fields along $c$.

\begin{theorem}[Second variation formula]
    \[
    \begin{aligned}
        \frac{d^2}{ds^2} L(c_s)|_{s=0} = \frac{1}{\|c_0`(0)\|}\{[\langle \nabla^c_{\frac{\partial}{\partial s}}, T \rangle]_{t=0}^{t=\ell}+I_{c_0}\langle X^{\perp}, X^{\perp}\rangle \}
    \end{aligned}
    \]
    All terms on the RHS are evaluated at $s=0$.
    \label{thm:13.1}
\end{theorem}
\begin{proof}
    \[
    \begin{aligned}
        \frac{d^2}{ds^2} L(c_s)|_{s=0} = \int_0^{\ell} \frac{d^2}{ds^2} \|T\| dt.
    \end{aligned}
    \]
    Since $[\frac{\partial}{\partial s}, \frac{\partial}{\partial t}]=0$:
    \[
    \begin{aligned}
        \frac{d}{ds} \|T\| = \frac{1}{2 \|T\|} \frac{\partial}{\partial s} \langle T, T \rangle = \frac{1}{\|T\|} \langle \nabla^c_{\frac{\partial}{\partial t}} T,T \rangle = \frac{1}{\|T\|} \langle \nabla_{\frac{\partial}{\partial t}}^c X, T \rangle.
    \end{aligned}
    \]
    Thus 
    \[
    \begin{aligned}
        \frac{\partial^2}{\partial s^2} \|T\| =& \frac{1}{\|T\|} \frac{\partial}{\partial s} \langle \nabla^c_{\frac{\partial}{\partial s}} X, T \rangle - \frac{1}{\|T\|^2} \langle \nabla^c_{\frac{\partial}{\partial s}} X, T \rangle \frac{\partial}{\partial s} \|T\| \\
        =&\frac{1}{\|T\|} \left( \langle \nabla^c_{\frac{\partial}{\partial s}} \nabla^c_{\frac{\partial}{\partial t}} X, T \rangle + \langle \nabla^c_{\frac{\partial}{\partial t}} X, \nabla^c_{\frac{\partial}{\partial s}} T \rangle \right) \\
        &-\frac{1}{\|T\|^3} \langle \nabla^c_{\frac{\partial}{\partial t}} X,T \rangle^2 \\
        =& \frac{1}{\|T\|} ( \langle \nabla^c_{\frac{\partial}{\partial t}} \nabla^c_{\frac{\partial}{\partial s}} X, T \rangle + \langle R(X,T)X, T \rangle \\
        &+ \langle \nabla^c_{\frac{\partial}{\partial t}} X, \nabla^c_{\frac{\partial}{\partial t}} X \rangle )- \frac{1}{\|T\|^3} \langle \nabla^c_{\frac{\partial}{\partial t}} X, T \rangle^2.
    \end{aligned}
    \]
    At $s=0$, $\nabla^c_{\frac{\partial}{\partial t}} X |_{s=0}= \nabla_{c_0} X$. Since $c_0$ is a geodesic,
    $\|T\|_{s=0} = \|c_0`(0)\|$ and $\nabla_{c_0}T=0$. Thus:
    \[
    \begin{aligned}
        \frac{d^2}{ds^2} \|T\| |_{s=0} =& \frac{1}{\|c_0`(0)\|} (\frac{d}{dt} \langle \nabla^c_{\frac{\partial}{\partial t}} X, T \rangle - \langle R(X,T)T,X \rangle \\
        &+ \langle \nabla_{c_0} X, \nabla_{c_0}X \rangle - \langle \nabla_{c_0}X, \frac{T}{\|T\|} \rangle^2).
    \end{aligned}
    \]
    Covariant differentiation of $X = X^{\perp} + \langle X, \frac{T}{\|T\|} \rangle \frac{T}{\|T\|}$ yields:
    $\nabla_{c_0} X = \nabla_{c_0} X^{\perp} + \langle \nabla_{c_0} X, \frac{T}{\|T\|}\rangle \frac{T}{\|T\|}$.
    Since $\langle X^{\perp}, T \rangle = 0 \implies \langle \nabla_{c_0} X^{\perp},T \rangle =0:$
    \[
    \begin{aligned}
        \langle \nabla_{c_0} X, \nabla_{c_0} X \rangle - \langle \nabla_{c_0}X, \frac{T}{\|T\|} \rangle^2 = \langle \nabla_{c_0} X^{\perp}, \nabla_{c_0} X^{\perp} \rangle.
    \end{aligned}
    \]
    Thus:
    \[
    \begin{aligned}
        \frac{d^2}{ds^2} \|T\| |_{s=0} =& \frac{1}{\|c_0`(0)\|}(\frac{d}{dt} \langle \nabla^c_{\frac{\partial}{\partial s}} X, T \rangle - \langle R(X,T)T,X \rangle \\
        &+\langle \nabla_{c_0} X^{\perp}, \nabla_{c_0} X^{\perp} \rangle ).
    \end{aligned}
    \]
    Integrating with respect to $t$ from $0$ to $\ell$ yields the result.
\end{proof}

From a Riemannian manifold $M$, the diameter is defined as:
\[
\begin{aligned}
    diam(M) := \sup_{p,q \in M} d(p,q) \qquad (\leq + \infty)
\end{aligned}
\]

\begin{theorem}[Bonnet-Myers theorem]
    Let $M$ be a complete n-dimensional Riemannian manifold ($n \geq 2$) such that
    $Ric \geq (n-1)k$ for the same $k>0$.
    Then $M$ is compact and $diam(M) \leq \frac{\pi}{\sqrt{k}}$
\end{theorem}

\begin{problem}[Final problem]
    Let $n \geq 2$, $k>0$, and $M:= \mathbb{S}^n (\frac{1}{\sqrt{k}})$ be the $n$-dimensional sphere of radius $\frac{1}{\sqrt{k}}$.
    Answer the following:
    \begin{enumerate}
        \item Show that the sectional curvature of $M$ is $k$. \\
        (By \autoref{prop:9.7}, this implies $Ric = (n-1)k$.)
        \item Show that $diam(M) = \frac{\pi}{\sqrt{k}}$
    \end{enumerate}
    \label{prob:20}
\end{problem}
Problem \autoref{prob:20} shows \autoref{thm:13.1} is optimal. It is known that under the assumption of the Bonnet-Myers theorem if
$diam(M) = \frac{\pi}{\sqrt{k}}$, then $M$ is isometric to $\mathbb{S}^n (\frac{1}{\sqrt{k}})$.